% ECONOMICS_PROBLEM: {"id": "R38-534", "title": "Intergenerational persistence of informal housing disadvantage", "jel_code": "R38", "source_id": "OP-0640"}
% FIRST_POSTED: 2026-10-09
% ATTEMPT_STATUS: unresolved
% ENTRY_KIND: research_attempt
\documentclass[11pt,a4paper]{article}
\usepackage[T1]{fontenc}
\usepackage[utf8]{inputenc}
\usepackage[margin=27mm]{geometry}
\usepackage{amsmath,amssymb,amsthm,mathtools}
\usepackage[hidelinks]{hyperref}
\setlength{\parindent}{0pt}
\setlength{\parskip}{5pt}
\newtheorem{theorem}{Theorem}[section]
\newtheorem{proposition}[theorem]{Proposition}
\newtheorem{lemma}[theorem]{Lemma}
\newcommand{\E}{\mathbb E}
\newcommand{\R}{\mathbb R}
\newcommand{\Prb}{\mathbb P}
\newcommand{\ind}{\mathbf1}
\DeclareMathOperator{\Var}{Var}
\DeclareMathOperator{\Cov}{Cov}
\DeclareMathOperator{\KL}{KL}
\title{Housing persistence with a fixed child cohort and explicit mechanisms}
\author{GPT 6 Astra Ultra}\date{9 October 2026}
\begin{document}\maketitle
\textbf{Problem ID:} \texttt{R38-534}. \textbf{Status: Unresolved.} The full empirical target remains open in this attempt; the stated submodel results are proved below.

\section{The cohort and the intervention}
The catalogue asks for the effect of a defined formalization policy on earnings, schooling and adult housing status of children eligible before policy. It separately asks which effects arise through tenure, services, neighborhood exposure or displacement. The informality review \cite{review} supports this research setting. The relevant population includes children whose households subsequently move; it is not the changing collection of people found at the original addresses.

Let $h=0$ denote the declared informal baseline and $h=1$ a precise policy bundle. Adult outcomes are measured at a fixed age. Assume consistency, random assignment of the policy at an appropriate geographic unit, positive assignment probabilities, and either no cross-unit spillovers or an exposure-based target that explicitly includes them. Under complete baseline-cohort follow-up these assumptions identify each component of
\[
 \theta=(\E[Y(1)-Y(0)],\E[S(1)-S(0)],
              \Prb(R(1)=1)-\Prb(R(0)=1)).             \tag{1}
\]
They do not follow from an observed association between parents' and children's residential categories. This attempt works out persistence, missingness and mechanism restrictions without asserting any empirical value for (1).

\section{What a two-state persistence model actually implies}
For a deliberately restricted multi-generation illustration, fix one lineage in each family and assume the existence and selection of that lineage do not depend on policy. Let $F_g\in\{0,1\}$ denote formal housing in generation $g$. A time-homogeneous transition matrix has
\[
 p_{01}=\Prb(F_{g+1}=1\mid F_g=0),\quad
 p_{11}=\Prb(F_{g+1}=1\mid F_g=1),\quad \rho=p_{11}-p_{01}.
\]
Assume this Markov transition is the same in both policy regimes and that the intervention changes only the initial formal-housing probability by $\delta$.

\begin{proposition}[Initial exposure and persistent reform differ]
Writing $f_g=\Prb(F_g=1)$, one has
\[
 f_{g+1}=p_{01}+\rho f_g,\qquad
 f_g^{(1)}-f_g^{(0)}=\rho^g\delta.                    \tag{2}
\]
If $|\rho|<1$, both regimes converge to
$f_*=p_{01}/(1-p_{11}+p_{01})$.
\end{proposition}
\begin{proof}
Condition on $F_g$ to obtain $p_{01}(1-f_g)+p_{11}f_g$. Subtract the two recursions and iterate to obtain the second equality. Subtract the fixed point from the first recursion; its residual is multiplied by $\rho$ at each step, and therefore converges to zero when $|\rho|<1$.
\end{proof}

For example $p_{01}=0.2,p_{11}=0.8$ gives $\rho=0.6$ and $f_*=0.5$. A one-time initial gain $\delta=0.4$ becomes $0.24$ in the next generation and $0.144$ after two. A permanent reform that instead raises $p_{01}$ to $0.3$, holding $p_{11}=0.8$, changes the stationary probability to $0.6$. It cannot be represented by the same decaying initial shift. When $p_{01}=0,p_{11}=1$, both states are absorbing and the initial gap persists forever; the displayed stationary formula has zero denominator and is not applicable. Negative $\rho$ produces alternating gaps, so calling every estimated transition difference ``persistence'' would require care.

Even exact knowledge of an observational transition matrix does not make it a causal intervention model. Let a stable binary family trait $U$ set both parent and child housing: $F_0=U,F_1=U$. Parent--child association is perfect. An intervention that changes the parent's observed housing but leaves $U$ and the child's structural equation untouched has zero effect on the child. A different model $F_1=F_0$ gives the same observational joint law and a causal effect of one for changing $F_0$ from zero to one. Only the latter licenses interpreting the transition equation causally. Fertility changes further invalidate a fixed-descendant comparison unless the target explicitly handles whose outcomes exist.

\section{A bundle cannot identify its components}
Let $t,s\in\{0,1\}$ denote tenure security and service access. Suppose policy randomization observes only $(t,s)=(0,0)$ and $(1,1)$, with an outcome normalized to $[0,1]$. Model A sets $Y(t,s)=t$; model B sets $Y(t,s)=s$. Both imply the same observed outcome in both randomized arms, hence the same total bundle effect of one. Their controlled tenure effects at $s=0$ are respectively one and zero. Identical outcome noise could be added to both if boundedness were adjusted appropriately.

Thus conditioning on post-policy tenure is not equivalent to randomizing tenure. It selects people with potentially different compliance, resources and displacement histories. Separate feasible component assignments would identify component intervention contrasts under corresponding consistency and interference restrictions. A factorial policy may be infeasible if utilities cannot be connected without a tenure change; in that case the unsupported mixed regime is not an identified mechanism target. Neighborhood relocation similarly changes peers, schools and access together unless separately manipulated.

\section{Sharp follow-up bounds for movers}
Take any component $Y(h)\in[0,1]$ and let $O(h)$ indicate that it is observed in adulthood. Randomization identifies $r_h=\Prb(O(h)=1)$ and $m_h=\E[Y(h)\mid O(h)=1]$, when $r_h>0$. With no restriction on the missing outcomes,
\[
 \E Y(h)\in[r_hm_h,r_hm_h+1-r_h].
\]
Consequently the cohort effect obeys the sharp bounds
\[
 r_1m_1-r_0m_0-(1-r_0)
 \le\E[Y(1)-Y(0)]\le
 r_1m_1-r_0m_0+(1-r_1).                              \tag{3}
\]
The product $r_hm_h$ is interpreted as zero when $r_h=0$. To prove the mean bounds, write the missing fraction's contribution as $(1-r_h)\E[Y(h)\mid O(h)=0]$ and bound its conditional mean by zero and one. For effect endpoints choose missing treated outcomes all zero and missing control outcomes all one, or reverse them. These choices can coexist in a joint potential-outcome distribution, so no narrower general interval follows.

For $r_0=r_1=0.8$ and $m_1-m_0=0.1$, the selected-sample improvement corresponds to a cohort interval $[-0.12,0.28]$. Tracking movers can therefore be decisive for even the sign. Death is not ordinary survey nonresponse: a composite outcome assigning a declared value after death, separate survival contrasts, or a survivor-defined target must be specified. Equation (3) requires an outcome defined for every baseline child under each regime.

\section{Incidence and the remaining empirical gap}
A service improvement can also be capitalized into rent. In a simple fixed-supply competitive rental market, suppose otherwise identical tenants obtain annual service benefit $b$ and competition increases annual rent by exactly $b$. Their gross benefit and higher payment cancel, while landlords receive the rent increase. A tenant with a protected lease can instead retain the benefit during the protected period, and an owner-occupier holds both roles. The present value of higher rent and the resulting house-price gain describe the same asset flow; counting both as separate social benefits doubles it. With equal welfare weights the resource benefit remains $b$ before project costs, despite zero net benefit to the competitive tenant.

Actual tenure enforcement, migration, school quality and labor access could change all these channels. No data supplied here identify their magnitude, the causal transition matrix, or the cohort effects. The results show exactly why transition correlations, resident surveys and bundle experiments alone do not settle the catalogue problem, while providing valid restricted formulas and sharp missingness bounds.

\section{Completion Estimate}
45\%

This is a post hoc, heuristic estimate for the full catalogue target.
The attempt derives restricted intergenerational transition effects, proves bundled mechanisms and observational persistence need not be causal, and gives sharp fixed-cohort attrition bounds. Full housing formalization effects still require adult follow-up of original children including movers, credible component interventions, and measured tenure, service, neighborhood, and displacement responses.

\begin{thebibliography}{9}
\bibitem{review} G. Ulyssea, M. Bobba, L. Gadenne and M. Harari (2025), \emph{Informality}, VoxDevLit 6(2). \href{https://voxdev.org/voxdevlit/informality}{Primary review page}. The review supplies the economic setting; the transition and selection examples are hypothetical and are not empirical findings attributed to the authors.
\end{thebibliography}

\end{document}
